\section{Desigualdades}

\begin{thm:desigualdade de markov}[da Desigualdade de Markov]
Se $X$ é uma variável aleatória não negativa, então para todo $a > 0$
\[
    P(X \ge a) \le \frac{\ev{X}}{a}
\]
\end{thm:desigualdade de markov}
\begin{proof}
Seja $I$ a função indicadora tal que
\[
    I(x) = \begin{cases}
        1,&\text{se } x \ge a \\
        0,&\text{caso contrário.}
    \end{cases}
\]
Com isso,
\[
    I \le \frac{X}{a},
\]
e tomando a esperança de ambos os lados
\[
    \ev{I} \le \frac{\ev{X}}{a}
\]
Já que
\[
    \ev{I} = 1 \cdot P(I = 1) = P(X \ge a),
\]
então
\[
    P(X \ge a) \le \frac{\ev{X}}{a}
\]
\end{proof}

\begin{cor:desigualdade de chebyshev}
\label{cor:desigualdade de chebyshev}
Se $X$ é uma variável aleatória com média $\mu$ e variância $\sigma^2$, então
para todo $k > 0$
\[
    P(|X - \mu| \ge k) \le \frac{\sigma^2}{k^2}
\]
\end{cor:desigualdade de chebyshev}
\begin{proof}
Pela desigualdade de Markov,
\[
    P((X - \mu)^2 \ge k^2) \le \frac{\ev{(X - \mu)^2}}{k^2}
\]
mas como $(X - \mu)^2 \ge k^2$ se, e somente se, $|X - \mu| \ge k$, e pela definição
de variância, temos que
\[
    P(|X - \mu| \ge k) \le \frac{\sigma^2}{k^2}
\]
\end{proof}

\begin{prop:variancia 0 e igual media}
Se $X$ é uma variável aleatória com $\var{X} = 0$, então $X = \ev{X}$.
\end{prop:variancia 0 e igual media}
\begin{proof}
Pela desigualde de Chebyshev,
\[
    P\left(|X - \mu| > \frac{1}{n}\right) \le 0
\]
e por não negatividade de probabilidades,
\[
    P\left(|X - \mu| > \frac{1}{n}\right) = 0
\]
para todo $n \in \real$.
Sabemos então que
\[
    \lim_{n \to \infty} P\left(|X - \mu| > \frac{1}{n}\right) = 0,
\]
no entanto,
\[
    \lim_{n \to \infty} P\left(|X - \mu| > \frac{1}{n}\right) =
    P\left[\lim_{n \to \infty} \left(|X - \mu| > \frac{1}{n}\right)\right] =
    P(|X - \mu| > 0),
\]
de modo que
\[
    P(|X - \mu| > 0) = 0.
\]
Portanto,
\[
    P(X = \mu) = 1 \Longrightarrow X = \ev{X}.
\]
\end{proof}

\begin{thm:desigualdade de jansen}
Se $X$ é uma variável aleatória e $g$ é uma função convexa, então
\[
    \ev{g(X)} \ge g(\ev{X})
\]
\end{thm:desigualdade de jansen}

\begin{cor:variancia positiva}
A variância de uma variável aleatória é sempre não negativa.
\end{cor:variancia positiva}
\begin{proof}
A função $g(x) = x^2$ é convexa, e pela Desigualdade de Jansen segue que
\[
    \ev{X^2} \ge \ev{X}^2.
\]
Uma vez que $\var{X} = \ev{X^2} - \ev{X}^2$, então $\var{X} \ge 0$.
\end{proof}

% \begin{thm:lei fraca dos grandes numeros}[Lei Fraca dos Grandes Números]
% Sejam $X_1, \ldots, X_n$ variáveis aleatória iid, logo $\ev{X_i} = \mu$ e
% $\var{X_i} = \sigma^2$ para todo $i \in [n]$.
% Para qualquer $\epsilon > 0$,
% \[
%     \lim_{n \to \infty} P\left(\left| \frac{\sum_{i=1}^n X_i}{n} - \mu \right| \ge \epsilon \right) = 0
% \]
% \end{thm:lei fraca dos grandes numeros}
% \begin{proof}
% Pelos Teoremas \ref{thm:propriedades esperanca} e \ref{prop:propriedades variancia},
% \[
%     \ev{\frac{\sum_{i=1}^n X_i}{n}} = \frac{n\mu}{n} = \mu \quad \text{e} \quad
%     \var{\frac{\õsum_{i=1}^n X_i}{n}} = \frac{n\sigma^2}{n^2} = \frac{\sigma^2}{n}.
% \]
% Pela desigualde de Chebyshev,
% \[
%     P\left(\left| \frac{\sum_{i=1}^n X_i}{n} - \mu \right| \ge \epsilon \right) \le \frac{\sigma^2}{n\epsilon}.
% \]
% Logo, temos
% \begin{gather*}
%     0 \le
%     \lim_{n \to \infty} P\left(\left| \frac{\sum_{i=1}^n X_i}{n} - \mu \right| \ge \epsilon \right) \le
%     \lim_{n \to \infty} \frac{\sigma^2}{n\epsilon} \\
%     0 \le
%     \lim_{n \to \infty} P\left(\left| \frac{\sum_{i=1}^n X_i}{n} - \mu \right| \ge \epsilon \right) \le
%     0
% \end{gather*}
% e, pelo Teorema do Confronto do cálculo de limties,
% \[
%     \lim_{n \to \infty} P\left(\left| \frac{\sum_{i=1}^n X_i}{n} - \mu \right| \ge \epsilon \right) = 0
% \]
% \end{proof}
